% Part 4: advanced usage examples. Divider plus 10 frames, session 3:40-3:55.
\section{Part 4. Advanced usage examples}\subsection{}

% Part 4 divider.
\begin{transitionframe}
Part 4. Advanced usage examples
\end{transitionframe}

% Framing: each demo adds an axis Part 2 could not show.
\begin{frame}{Each of these adds something the project could not show}
\begin{center}
\tabfont
\begin{tabular}{lll}
\toprule
Demo & Axis it adds & Where Part 2 stopped \\
\midrule
TravelTab + PetTab & The pipeline becomes \emph{software} & We shipped a report, not a system \\
GBrain    & Retrieval that is \emph{indexed} & Step 3 attached sample docs by hand \\
NanoClaw  & The agent has an \emph{address}  & You were always at the terminal \\
\bottomrule
\end{tabular}
\end{center}
\vs
Nothing here is a new tool. It is the same agent, moved one step along one axis.
\end{frame}

% TravelTab, 1 of 3: the background, as an annoyance.
\begin{frame}{TravelTab: personal-database management}
\begin{wideitemize}
  \item Background: I travel between New Haven and Philadelphia constantly, and Amtrak is
        the only workable way to do it. The app is painful, it does not talk to my
        calendar, \ybb{the ticket arrives once by email and sometimes it is gone}, and
        the pricing is dynamic.
  \item So I built \ybb{TravelTab}, a personal Amtrak ticket tracker:
  \begin{enumerate}\small
      \item \ylbb{Ingests tickets automatically from my email.}
      \item \ylbb{Exports structured records into a database}: one row per
            passenger per ticket.
      \item \ylbb{Syncs into my calendar}, so a trip is something I can actually
            manage.
      \item \ylbb{Fare watch} that checks prices and alerts me.
  \end{enumerate}
  \item I can access it from anywhere, on any device, with a web browser.
\end{wideitemize}
\end{frame}

% TravelTab, 2 of 3: the tickets dashboard.
\begin{frame}{TravelTab: personal-database management}
\begin{center}
\includegraphics[height=0.8\textheight]{traveltab_tickets1.png}
\end{center}
\vspace{-0.2em}
Every row here arrived as \ybb{an email}. Messy documents in, structured
records out.
\end{frame}

% TravelTab, 3 of 3: the fare watch, which runs unattended.
\begin{frame}{And it keeps running when you are not watching}
\begin{center}
\includegraphics[height=0.8\textheight]{traveltab_fares.png}
\end{center}
\vspace{-0.2em}
\ybb{Fare watchlist} that checks prices on a schedule and alerts me on a threshold.
\end{frame}

% PetTab, 1 of 3: same stack, another domain.
\begin{frame}{PetTab: another domain, same stack}
\begin{columns}[T]
\begin{column}{0.5\textwidth}
\begin{center}

\includegraphics[height=0.5\textheight]{petab_winnie.jpeg}
\end{center}
\end{column}
\begin{column}{0.5\textwidth}
\begin{center}

\includegraphics[height=0.5\textheight]{petab_dory.jpg}
\end{center}
\end{column}
\end{columns}
\vs
\begin{itemize}
    \item Background: I have two cats, Winnie and Dory, and I want to track their health and bahavior. Existing tools are too complicated, and I want full control over the data.
    \item So I built \ybb{PetTab}, a pet-health logger built on TravelTab's stack.
\end{itemize}
\end{frame}

% PetTab, 2 of 3: the data views.
\begin{frame}{PetTab: visualizing the data}
\begin{center}
\includegraphics[width=0.9\textwidth]{pettab_list.png}
\end{center}
\end{frame}

% PetTab, 3 of 3: an agent inside the product, on a read-only snapshot.
\begin{frame}{PetTab: a chat agent inside the product}
\begin{columns}[T]
\begin{column}{0.55\textwidth}
\begin{center}
\includegraphics[width=\textwidth]{pettab_chat.png}
\end{center}
\end{column}
\begin{column}{0.45\textwidth}
\vs
\begin{itemize}
    \item An agent is embedded within the product.
    \item PetTab's Insights chat runs the \file{claude} CLI over a
        \ybb{read-only snapshot} of the database, so that it can explore the database and give me suggestions or help me manage the reminders.
\end{itemize}
\end{column}
\end{columns}
\end{frame}

% GBrain: indexed retrieval over a corpus you own.
\begin{frame}{GBrain: personal-knowledge brain}
\begin{columns}[T]
\begin{column}{0.5\textwidth}

\includegraphics[width=\textwidth]{gbrain1.png}
\includegraphics[width=\textwidth]{gbrain2.png}
\includegraphics[width=\textwidth]{gbrain3.png}
\end{column}
\begin{column}{0.5\textwidth}
\vs
\begin{itemize}
    \item \ybb{GBrain} indexes a corpus you own -- meeting notes, summaries, Notion --
        and exposes it to the agent as a tool it can query.
    \item The key idea: it (1) searches the corpus, (2) reads the hits for you, and (3) writes the answer.
    \item {\small \yb{\url{github.com/garrytan/gbrain}}}
\end{itemize}
\end{column}
\end{columns}
\end{frame}

% NanoClaw: an agent you reach rather than launch.
\begin{frame}{An agent with an address}
\begin{columns}[T]
\begin{column}{0.5\textwidth}
\begin{center}
\includegraphics[height=0.75\textheight]{nanoclaw.png}
\end{center}
\end{column}
\begin{column}{0.5\textwidth}
\begin{itemize}
    \item \ybb{NanoClaw}: Claude in a Linux container, reached from \yb{WhatsApp,
    Slack, Telegram or email}.
    \item Every minute so far had you at a terminal, driving. Same tool -- reached
    by sending it a message.
    \item A personal-scale project: one process and a handful of files, running in a Linux container.
    \item {\small\yb{\url{github.com/nanocoai/nanoclaw}}}
\end{itemize}
\end{column}
\end{columns}
\end{frame}

% Where to go next, and the repo link.
\begin{frame}{Where to go next}
\begin{itemize}
    \item Try to use it in your own projects and daily life \yob{as much as possible}. 
    \item Brainstorm your ideas, and then \yob{build them}.
\end{itemize}

Other useful resources:
\begin{itemize}
    \item \yb{Agent Skills with Anthropic} -- DeepLearning.AI short course
    \item \yb{Finance skill plugins} by Anthropic 
    \item A collection of resources on AI tools: \yb{\url{www.xuganchen.com/public-goods}}
    \item Alternatively, ask AI to build it for you.
\end{itemize}

\vs
Course materials and every prompt run today:\\ \yb{\url{github.com/xuganchen/BootCamp-Introduction-to-Claude-Code}}

\end{frame}
